% Part: first-order-logic
% Chapter: completeness
% Section: henkin-expansion

\documentclass[../../../include/open-logic-section]{subfiles}

\begin{document}

\olfileid[id]{fol}{com}{hen}

\olsection{Perluasan Henkin}

\begin{explain}
Salah satu tantangan dalam membuktikan teorema kelengkapan adalah bahwa
model yang kita bangun dari himpunan lengkap dan konsisten~$\Gamma$ harus
membuat semua !!{formula}s berkuantor dalam~$\Gamma$ benar. Untuk
menjaminnya, kita menggunakan siasat yang berasal dari Leon Henkin. Pada
intinya, siasat tersebut terdiri atas perluasan bahasa dengan tak hingga
banyak !!{constant}s serta penambahan, untuk setiap !!{formula} dengan
satu !!{variable} bebas $!A(x)$, suatu formula berbentuk
\iftag{prvEx}
      {$\lexists[x][!A(x)] \lif !A(c)$}
      {$\lnot\lforall[x][!A(x)] \lif \lnot !A(c)$},
dengan $c$ sebagai salah satu di antara !!{constant}s baru. Ketika kita membangun
!!{structure} yang memenuhi~$\Gamma$, cara ini menjamin bahwa setiap
\iftag{prvEx}
{kalimat eksistensial yang benar memiliki saksi}
{kalimat universal yang salah memiliki contoh penyangkal}
di antara konstanta-konstanta baru.
\end{explain}

\begin{prop}
\ollabel{prop:lang-exp}
Jika $\Gamma$ konsisten dalam $\Lang L$ dan $\Lang L'$ diperoleh dari
$\Lang L$ dengan menambahkan suatu himpunan yang !!a{denumerable}, dengan
anggota berupa !!{constant}s baru $\Obj d_0$, $\Obj d_1$, \dots, maka $\Gamma$
konsisten dalam~$\Lang L'$.
\end{prop}

\begin{defn}[Himpunan jenuh]
Suatu himpunan $\Gamma$ dari !!{formula}s dalam bahasa $\Lang {L}$
bersifat \emph{jenuh} jika dan hanya jika, untuk setiap
!!{formula}~$!A(x) \in \Frm[L]$ dengan satu !!{variable} bebas~$x$,
terdapat !!a{constant}~$c \in \Lang{L}$ sedemikian sehingga
\iftag{prvEx}
      {$\lexists[x][!A(x)] \lif !A(c) \in \Gamma$}
      {$\lnot\lforall[x][!A(x)] \lif \lnot !A(c) \in \Gamma$}.
\end{defn}

Definisi berikut akan digunakan dalam pembuktian teorema selanjutnya.

\begin{defn}
\ollabel{defn:henkin-exp}
Misalkan $\Lang L'$ seperti dalam \olref{prop:lang-exp}. Tetapkan suatu
enumerasi $!A_0(x_0)$, $!A_1(x_1)$, \dots dari semua
!!{formula}s~$!A_i(x_i)$ dalam~$\Lang L'$ yang mempunyai satu variabel
bebas ($x_i$). Kita mendefinisikan !!{sentence}s~$!D_n$ melalui
induksi pada~$n$.

Misalkan $c_0$ adalah !!{constant} pertama di antara $\Obj d_i$ yang
kita tambahkan ke~$\Lang{L}$ yang tidak muncul dalam~$!A_0(x_0)$.
Dengan mengandaikan bahwa $!D_0$, \dots,~$!D_{n-1}$ telah didefinisikan,
misalkan $c_n$ adalah yang pertama di antara !!{constant}s baru~$\Obj d_i$ yang tidak
muncul dalam $!D_0$, \dots,~$!D_{n-1}$ maupun dalam~$!A_n(x_n)$.

Sekarang, misalkan $!D_{n}$ adalah !!{formula} \iftag{prvEx}
{$\lexists[x_{n}][!A_{n}(x_{n})] \lif
  !A_{n}(c_{n})$}{$\lnot\lforall[x_{n}][!A_{n}(x_{n})] \lif \lnot
  !A_{n}(c_{n})$}.
\end{defn}

\begin{lem}
\ollabel{lem:henkin}
Setiap himpunan konsisten~$\Gamma$ dapat diperluas menjadi himpunan
konsisten yang jenuh~$\Gamma'$.
\end{lem}

\begin{proof}
Diberikan suatu himpunan kalimat yang konsisten~$\Gamma$ dalam
bahasa~$\Lang{L}$, perluas bahasa tersebut dengan menambahkan suatu
himpunan yang !!a{denumerable}, dengan anggota berupa !!{constant}s baru,
sehingga terbentuk~$\Lang{L'}$.
Menurut \olref{prop:lang-exp}, $\Gamma$ tetap konsisten dalam bahasa yang
lebih kaya ini. Selanjutnya, misalkan $!D_i$ seperti dalam
\olref{defn:henkin-exp}. Misalkan
\begin{align*}
\Gamma_0 & = \Gamma \\
\Gamma_{n+1} & = \Gamma_n \cup \{!D_n \}
\end{align*}
yakni, $\Gamma_{n+1} = \Gamma \cup \{ !D_0, \dots, !D_n \}$, dan
misalkan $\Gamma' = \bigcup_{n} \Gamma_n$. Jelas bahwa $\Gamma'$ jenuh.

Jika $\Gamma'$ tak konsisten, maka untuk suatu $n$, $\Gamma_n$ tak
konsisten (Latihan: jelaskan mengapa). Jadi, untuk menunjukkan bahwa
$\Gamma'$ konsisten, cukup ditunjukkan, melalui induksi pada~$n$, bahwa
setiap himpunan~$\Gamma_n$ konsisten.

Basis induksinya hanyalah klaim bahwa $\Gamma_0 = \Gamma$ konsisten,
yang merupakan hipotesis teorema. Untuk langkah induksi, andaikan bahwa
$\Gamma_{n}$ konsisten, tetapi $\Gamma_{n+1} = \Gamma_n \cup \{!D_n\}$
tak konsisten. Ingat bahwa $!D_n$~adalah
\iftag{prvEx}
{$\lexists[x_{n}][!A_{n}(x_n)] \lif !A_{n}(c_{n})$}
{$\lnot\lforall[x_{n}][!A_{n}(x_n)] \lif \lnot !A_{n}(c_{n})$},
dengan $!A_n(x_n)$ sebagai !!a{formula} dalam $\Lang{L'}$ yang hanya
memuat variabel~$x_n$ secara bebas. Berdasarkan cara kita memilih~$c_n$
(lihat \olref{defn:henkin-exp}), $c_n$ tidak muncul dalam~$!A_n(x_n)$
maupun dalam~$\Gamma_n$.

Jika $\Gamma_n \cup \{!D_n\}$ tak konsisten, maka $\Gamma_n
\Proves \lnot !D_n$, sehingga kedua hal berikut berlaku:
\iftag{prvEx}{
\[
\Gamma_n \Proves \lexists[x_n][!A_n(x_n)]
\qquad
\Gamma_n \Proves \lnot !A_n(c_n)
\]}{
\[
\Gamma_n \Proves \lnot\lforall[x_n][!A_n(x_n)]
\qquad
\Gamma_n \Proves !A_n(c_n)
\]}
Karena $c_n$ tidak muncul dalam
$\Gamma_n$ maupun dalam~$!A_n(x_n)$,
\tagrefs{prfAX/{fol:axd:qpr:thm:strong-generalization},prfSC/{fol:seq:qpr:thm:strong-generalization},prfND/{fol:ntd:qpr:thm:strong-generalization},prfTab/{fol:tab:qpr:thm:strong-generalization}}
berlaku. Dari \iftag{prvEx}
{$\Gamma_n \Proves \lnot !A_n(c_n)$}
{$\Gamma_n \Proves !A_n(c_n)$},
kita memperoleh
\iftag{prvEx}
{$\Gamma_n \Proves \lforall[x_n][\lnot !A_n(x_n)]$}
{$\Gamma_n \Proves \lforall[x_n][!A_n(x_n)]$}.
Dengan demikian, kita memperoleh sekaligus
\iftag{prvEx}
{$\Gamma_n \Proves \lexists[x_n][!A_n(x_n)]$ dan
$\Gamma_n \Proves \lforall[x_n][\lnot !A_n(x_n)]$}
{$\Gamma_n \Proves \lnot\lforall[x_n][!A_n(x_n)]$ dan
$\Gamma_n \Proves \lforall[x_n][!A_n(x_n)]$},
sehingga $\Gamma_n$ sendiri tak konsisten.
\iftag{prvEx}
{(Perhatikan bahwa
\iftag{prvAll}
{$\lforall[x_n][\lnot !A_n(x_n)] \Proves
  \lnot\lexists[x_n][!A_n(x_n)]$}
{$\lforall[x_n][\lnot !A_n(x_n)]$ didefinisikan sebagai
  $\lnot\lexists[x_n][\lnot\lnot !A_n(x_n)]$ dan
  $\lnot\lexists[x_n][\lnot\lnot !A_n(x_n)] \Proves
  \lnot\lexists[x_n][!A_n(x_n)]$}.)}{}
Kontradiksi: $\Gamma_n$ seharusnya konsisten. Jadi,
$\Gamma_n \cup \{ !D_n\}$ konsisten.
\end{proof}

\begin{explain}
Sekarang akan kita tunjukkan bahwa himpunan yang \emph{lengkap},
konsisten, dan jenuh memiliki sifat bahwa \iftag{prvAll}{himpunan itu
  memuat suatu !!{sentence} berkuantor universal jika dan hanya jika
  himpunan itu memuat semua instansnya}{}\iftag{defAll,defEx}{}{ dan }
\iftag{prvEx}{himpunan itu memuat suatu !!{sentence} berkuantor
  eksistensial jika dan hanya jika himpunan itu memuat sekurang-kurangnya
  satu instansnya}{}. Kita akan menggunakan sifat ini untuk menunjukkan
bahwa !!{structure} yang akan kita hasilkan dari suatu himpunan lengkap,
konsisten, dan jenuh membuat semua kalimat berkuantornya benar.
\end{explain}

\begin{prop}\ollabel{prop:saturated-instances}
Andaikan $\Gamma$ lengkap, konsisten, dan jenuh.
\begin{tagenumerate}{prvEx,prvAll}
\tagitem{prvEx}{$\lexists[x][!A(x)] \in \Gamma$ jika dan hanya jika
  $!A(t) \in \Gamma$ untuk sekurang-kurangnya satu suku tertutup~$t$.}{}
\tagitem{prvAll}{$\lforall[x][!A(x)] \in \Gamma$ jika dan hanya jika
  $!A(t) \in \Gamma$ untuk semua suku tertutup~$t$.}{}
\end{tagenumerate}
\end{prop}

\begin{proof}
  \begin{tagenumerate}{prvEx,prvAll}
  \tagitem{prvEx}{%
    \iftag{probEx}{Latihan.}{Pertama, andaikan bahwa
      $\lexists[x][!A(x)] \in \Gamma$. Karena $\Gamma$ jenuh,
      $(\lexists[x][!A(x)] \lif !A(c)) \in \Gamma$ untuk suatu
      !!{constant}~$c$. Menurut
      \tagrefs{prfAX/{fol:axd:ppr:prop:provability-lif},prfSC/{fol:seq:ppr:prop:provability-lif},prfND/{fol:ntd:ppr:prop:provability-lif},prfTab/{fol:tab:ppr:prop:provability-lif}}, butir~(1),
      dan \olref[ccs]{prop:ccs}\olref[ccs]{prop:ccs-prov-in}, $!A(c)
      \in \Gamma$.

    Untuk arah sebaliknya, kejenuhan tidak diperlukan: andaikan
    $!A(t) \in \Gamma$. Maka $\Gamma \Proves \lexists[x][!A(x)]$
    menurut
      \tagrefs{prfAX/{fol:axd:qpr:prop:provability-quantifiers},prfSC/{fol:seq:qpr:prop:provability-quantifiers},prfND/{fol:ntd:qpr:prop:provability-quantifiers},prfTab/{fol:tab:qpr:prop:provability-quantifiers}}, butir~(1). Menurut
    \olref[ccs]{prop:ccs}\olref[ccs]{prop:ccs-prov-in},
    $\lexists[x][!A(x)] \in \Gamma$.}}{}
  \tagitem{prvAll}{%
    \iftag{probAll}{Latihan.}{Andaikan bahwa $!A(t) \in \Gamma$ untuk
      semua suku tertutup~$t$. Demi memperoleh kontradiksi, andaikan
      $\lforall[x][!A(x)] \notin \Gamma$. Karena $\Gamma$ lengkap,
      $\lnot\lforall[x][!A(x)] \in \Gamma$. Karena jenuh,
      \iftag{prvEx}{$(\lexists[x][\lnot !A(x)] \lif \lnot !A(c)) \in
        \Gamma$}{$(\lnot\lforall[x][!A(x)] \lif \lnot !A(c)) \in
        \Gamma$} untuk suatu !!{constant}~$c$. Menurut asumsi, karena
      $c$ merupakan suku tertutup, $!A(c) \in \Gamma$. Namun, hal ini
      akan membuat $\Gamma$ tak konsisten. (Latihan: berikan
      !!{derivation} yang menunjukkan bahwa
      \[
      \lnot \lforall[x][!A(x)],
      \iftag{prvEx}{\lexists[x][\lnot !A(x)] \lif \lnot
        !A(c)}{\lnot\lforall[x][!A(x)] \lif \lnot
        !A(c)}, !A(c)
      \]
      tak konsisten.)

      Untuk arah sebaliknya, kita tidak memerlukan kejenuhan: andaikan
      $\lforall[x][!A(x)] \in \Gamma$. Maka $\Gamma \Proves !A(t)$
      menurut
      \tagrefs{prfAX/{fol:axd:qpr:prop:provability-quantifiers},prfSC/{fol:seq:qpr:prop:provability-quantifiers},prfND/{fol:ntd:qpr:prop:provability-quantifiers},prfTab/{fol:tab:qpr:prop:provability-quantifiers}},
      butir~(2). Kita memperoleh $!A(t) \in \Gamma$ menurut
      \olref[ccs]{prop:ccs}.}}{}
  \end{tagenumerate}
\end{proof}

\end{document}
